\documentclass[a4paper]{article}
\usepackage{amsfonts}
\usepackage{amsmath}
\usepackage{amssymb}
\usepackage{graphicx}     

\begin{document}
  
\noindent \large Auteur: Raj Shah \\
\noindent \large Studierichting: Informatica\\
\noindent \large Oefeningenreeks 3A nr. 2.6\\

\medskip

\normalsize

Bereken de afgeleide van $f(x) = -3x^2$ voor $a \in \left\{-2, 1, 5\right\}$\\

$f'(a) = \lim\limits_{h\rightarrow 0} \dfrac{f(a+h) - f(a)}{h} $\\
\begin{itemize}
	\item $a = -2$\\

$f'(-2) = \lim\limits_{h\rightarrow 0} \dfrac{-3(h-2)^2 - f(-2)}{h} $\\

	$ = \lim\limits_{h\rightarrow 0} \dfrac{-3(h^2 -4h + 4) + 12}{h}$\\
	
	$ = \lim\limits_{h\rightarrow 0} \dfrac{-3h^2 +12h}{h}$\\ 
	
	$ = \lim\limits_{h\rightarrow 0} (-3h + 12) = 12$\\

	\item $a = 1$\\

$f'(1) = \lim\limits_{h\rightarrow 0} \dfrac{-3(h+1)^2 - f(1)}{h} $\\

	$ = \lim\limits_{h\rightarrow 0} \dfrac{-3(h^2 + 2h +1) -3}{h}$\\
	
	$ = \lim\limits_{h\rightarrow 0} \dfrac{-3h^2 -6h}{h}$\\
	
	$ = \lim\limits_{h\rightarrow 0} (-3h - 6) = -6 $\\
	
	\item $a = 5$\\

$f'(5) = \lim\limits_{h\rightarrow 0} \dfrac{-3(h+5)^2 - f(5)}{h} $\\

	$ = \lim\limits_{h\rightarrow 0} \dfrac{-3(h^2 + 10h + 25) -75}{h}$\\
	
	$ = \lim\limits_{h\rightarrow 0} \dfrac{-3h^2 - 30h}{h}$\\
	
	$ = \lim\limits_{h\rightarrow 0} (-3h -30) = -30$\\

\end{itemize}

\begin{thebibliography}{999}
%\bibitem{a} Referentie
\end{thebibliography}
\end{document} 
